\chapter*{Résumé}
\phantomsection
\addcontentsline{toc}{chapter}{Résumé}

\textbf{Ce mémoire s'inscrit dans le cadre d'un Projet de Fin d'Études} réalisé à l'EMSI Rabat (filière Ingénierie Informatique et Réseaux) et enrichi par un stage au sein de l'agence \textbf{Agence de Développement du Digital (\#ADD)}. Il vise la conception et la réalisation d'une plateforme logicielle de \emph{services électroniques} destinée à rationaliser la relation entre les usagers et les démarches administratives ou organisationnelles portées en ligne. \textbf{L'objectif principal consiste à proposer une solution modulaire, sécurisée et évolutive}, capable d'intégrer progressivement de nouveaux services sans remettre en cause l'architecture centrale.

\textbf{La démarche adoptée repose sur une ingénierie logicielle structurée}, articulée autour d'une présentation de l'agence d'accueil, d'une analyse du contexte et de la problématique, d'un cahier des charges fonctionnel et non fonctionnel, d'une étude de l'existant, puis d'une analyse des besoins formalisée sous forme de règles métiers, de cas d'utilisation et de contraintes de qualité. La conception combine des vues UML (cas d'utilisation, classes, séquences, composants et déploiement) interprétées de manière descriptive, une architecture technique en couches (présentation, API, domaine, persistance) et un modèle de données relationnel normalisé. \textbf{Les choix techniques sont étayés par des références bibliographiques et documentaires.}

\textbf{L'implémentation s'appuie sur un socle full stack cohérent} : une API REST Laravel (PHP), une interface React (Vite), une base PostgreSQL, une authentification \emph{SPA} via Laravel Sanctum (cookies de session et protection CSRF), ainsi qu'une chaîne de déploiement documentée (variables d'environnement, migrations Artisan, conteneurs Docker optionnels). Les fonctionnalités couvrent la gestion des comptes et rôles, la configuration dynamique des services, le dépôt et le suivi des demandes (y compris parcours public), les notifications e-mail, et les espaces d'administration.

\textbf{Les phases de tests et de validation} combinent tests unitaires, tests d'intégration, tests fonctionnels et scénarios de charge ciblés. Le déploiement est envisagé selon un modèle conteneurisé, avec séparation des environnements, variables de configuration externalisées et supervision minimale. \textbf{Les résultats obtenus démontrent la faisabilité technique du dispositif} et ouvrent des perspectives d'enrichissement fonctionnel, d'interopérabilité et d'industrialisation.

\vspace{0.8cm}
\noindent\textbf{Mots-clés :} services électroniques, architecture REST, Laravel, React, PostgreSQL, Sanctum, sécurité, conteneurisation, tests, déploiement.

\cleardoublepage